\documentclass[10pt,a4paper]{amsart}
\usepackage[margin=19mm]{geometry}
\usepackage{amsmath,amssymb,mathtools}
\usepackage[scr=rsfs,cal=euler]{mathalfa}
\usepackage{xfrac}
\usepackage[unicode,hidelinks,bookmarks=false]{hyperref}
\newcommand{\ie}{i.e.\ }
\newcommand{\cf}{cf.\ }
\newcommand{\resp}{respectively\ }
\newcommand{\Subsection}{\subsection}
\providecommand{\nobreakdash}{\nobreak\hbox{-}}
\setlength{\emergencystretch}{2em}
\multlinegap0pt
\usepackage{enumerate}
\usepackage{amssymb}
\usepackage{mathtools}
\usepackage{tikz}
\newtheorem{thm}{Theorem}
\title{Jacobi Hamiltonian Integrators: construction and applications}
\author{Ad\'erito Ara\'ujo, Gon\c{c}alo Inoc\^encio Oliveira, Jo\~ao Nuno Mestre}
\date{Source: arXiv:2601.20799v1 (CC BY 4.0)}
\begin{document}
\maketitle

\noindent\emph{Source and licence.} Jacobi Hamiltonian Integrators: construction and applications. Version arXiv:2601.20799v1 (CC BY 4.0), distributed by arXiv under the Creative Commons Attribution 4.0 International licence, http://creativecommons.org/licenses/by/4.0/, as recorded in the arXiv licence metadata for this exact version and shown on the abstract page https://arxiv.org/abs/2601.20799v1. Full licence notice: \texttt{licenses/arxiv-2601.20799-CC-BY-4.0.txt}.\par\smallskip

\noindent\textit{Editorial scope.} Counted unit: the article's thm homogeneous (source0.tex lines 150--152). The article's complete original proof of it is delivered with it (source0.tex lines 154--179, one \texttt{proof} environment of 26 lines). No mathematics is written, repaired or re-derived: the statement and the proof are the article's own lines, in the article's own notation, and no retained line is substituted or re-flowed.  No further source-local context is needed: every cross-reference of every retained line resolves inside the two delivered spans.  Nothing was omitted from any retained span, so the package carries no omission record.  No retained line contains a citation, so the wrapper carries no bibliography and no bibliography entry.  No retained line contains a figure environment, an image-inclusion command or drawing code, so no image is delivered and the package declares no asset dependency.  Editorial formatting only, in addition: an AMS \texttt{amsart} preamble that restates the article's own declarations and macros used by the retained lines, namely the environments \texttt{thm} on separate counters, together with the standard packages those lines need.  The retained environments print in this document as Theorem 1 in order of appearance, whereas the article prints the same items with the numbering of its own full document.  The complete native source of the article as distributed in the arXiv e-print for this exact version is bundled unchanged as \texttt{sources/193550/source0.tex}.
\begin{thm}\label{thm: S_i homogeneous}
The recursion defined in steps 1 and 2 preserves the homogeneous structure; that is, all generated $S_i$ are 1-homogeneous.
\end{thm}

\begin{proof}
    We proceed by induction on $k$. For $k=1$, we have $S_1(x,t) = H(x,t)$, which is clearly 1-homogeneous. Assume now that $S_1, \dots, S_k$ are all 1-homogeneous. Consider $S_s^{(k)} = \sum_{j=1}^k s^j S_j$; by the induction hypothesis, $S_s^{(k)}$ is also 1-homogeneous.  

We claim that $d(S_s^{(k)} \circ h_z) = T^*h_z \circ d S_s^{(k)}$. Indeed, for the $x$-components, one has
\[
    \frac{\partial S_s^{(k)}}{\partial x}(x, zt) = z \frac{\partial S_s^{(k)}}{\partial x}(x, t),
\]
and for the $t$-component, using the chain rule,
\[
    \frac{d}{dt}S_s^{(k)} (x,zt) = z \frac{\partial S^{(k)}_s}{\partial t}(x,zt). 
\]

On the other hand, $\frac{d}{dt}S_s^{(k)} (x,zt) = \frac{d}{dt}\left(zS_s^{(k)} (x,t) \right) = z \frac{\partial S^{(k)}_s}{\partial t}(x,t)$. So,
\[\frac{\partial S_s^{(k)}}{\partial t} (x,zt) =\frac{\partial S_s^{(k)}}{\partial t} (x,t).\]
        

Finally, $S_{k+1}$ is 1-homogeneous because both $\alpha$ and $dS_s^{(k)}$ are homogeneous:
\begin{align*}
    h_z^* S_{k+1}(x,t) &= \frac{1}{(k+1)!} \frac{d^k}{ds^k}\Big|_{s=0} \hat{H}\Big( \alpha(d_x S_s^{(k)} \circ h_z) \Big)\\
    &= \frac{1}{(k+1)!} \frac{d^k}{ds^k}\Big|_{s=0} \hat{H}\Big( \alpha(T^*h_z \circ d_x S_s^{(k)}) \Big)\\
    &= \frac{1}{(k+1)!} \frac{d^k}{ds^k}\Big|_{s=0} \hat{H}\Big( h_z \circ \alpha(d_x S_s^{(k)}) \Big)\\
    &= z \frac{1}{(k+1)!} \frac{d^k}{ds^k}\Big|_{s=0} \hat{H}\Big( \alpha(d_x S_s^{(k)}) \Big)\\
    &= z S_{k+1}(x,t), 
\end{align*}
which concludes the proof.
\end{proof} 

\end{document}
